% =====================================================================
% P1-15 -- First full manuscript draft (LaTeX source).
% Content: frozen theorem core T1-T6, C1-C6 (corrected C5), proved in
% P1_04/P1_05/P1_06 and independently audited in P1_14.  The candidate
% corollaries N1-N3 of P1-14 are NOT used here; they are collected in
% P1_15_OPTIONAL_COROLLARIES.md pending author approval.
% No novelty or priority claim is made anywhere in this draft.
% P1-17 (21 Oct 2026): audit corrections of P1_16_INDEPENDENT_AUDIT.md
% applied; see P1_17_CORRECTION_LOG.md. Core statements unchanged.
% Compilation: pdflatex P1_15_MANUSCRIPT_DRAFT.tex (twice).  Only
% standard packages (amsmath, amssymb, amsthm, geometry) are required.
% =====================================================================
\documentclass[11pt]{article}

\usepackage[a4paper,margin=2.6cm]{geometry}
\usepackage{amsmath,amssymb,amsthm}

\newtheorem{theorem}{Theorem}[section]
\newtheorem{lemma}[theorem]{Lemma}
\newtheorem{proposition}[theorem]{Proposition}
\newtheorem{corollary}[theorem]{Corollary}
\theoremstyle{definition}
\newtheorem{definition}[theorem]{Definition}
\newtheorem{example}[theorem]{Example}
\theoremstyle{remark}
\newtheorem*{remark}{Remark}

\newcommand{\Fq}{\mathbb{F}_q}
\newcommand{\Fp}{\mathbb{F}_p}
\newcommand{\Z}{\mathbb{Z}}
\newcommand{\Hull}{\operatorname{Hull}}
\newcommand{\lcm}{\operatorname{lcm}}
\newcommand{\tr}{\operatorname{tr}}
\newcommand{\E}{\mathbb{E}}

\title{The joint code--hull distribution of repeated-root constacyclic codes
over affine algebras via transfer matrices}
\author{[Author names and affiliations to be inserted]}
\date{Draft of 21 October 2026}

\begin{document}
\maketitle

\begin{abstract}
Let \(A \cong \prod_{s=1}^{N} K_s\) be a square-free affine algebra over
\(\Fq\), let \(P = p^{\varrho}\) be a power of the characteristic, and let
\(n = P n_0\) with \(\gcd(n_0,p)=1\). For compatible twists \(\lambda_s\) we
study the \(\lambda_s\)-constacyclic codes of length \(n\) over each
\(K_s = \mathbb{F}_{q^{m_s}}\) in the repeated-root regime
\(x^n - \lambda_s = (x^{n_0} - \mu_s)^P\), parametrized by exponent
assignments \(a_f \in \{0,\dots,P\}\) on the distinct irreducible factors
of the separable polynomial \(x^{n_0} - \mu_s\). Using the code-first
\(k\)-Galois pairing, we give the Galois dual and hull of such a code in
exponent form: on each orbit of the reciprocal involution \(\tau_s\), the
hull dimension is a sum of edge terms \(\min(P-\alpha_i, \alpha_{i-1})\).
Our main result is a joint enumerator
\(\mathcal{E}(u,z) = \sum_a u^{\dim_{\Fq} C(a)} z^{\dim_{\Fq} \Hull_k(C(a))}
= \prod_s \prod_O \tr\!\big(T^{(P)}_{w_O}(u,z)^{L_O}\big)\),
where \(T^{(P)}_w(u,z)_{a,b} = u^{w(P-b)} z^{w\min(P-b,a)}\) is a
\((P{+}1)\)-state transfer matrix attached to each \(\tau_s\)-orbit. We
derive the total count, the code-dimension distribution, the LCD count
\(2^{\#\mathrm{orbits}}\) (independent of \(P\)), characterizations and
counts of self-orthogonal and dual-containing codes, and exact mean and
variance formulas for the hull dimension, in which orbits of length one
require a separate term \(\nu_P = \lfloor P^2/4\rfloor/(P+1)\). At \(P=1\)
the enumerator specializes to the binary-state simple-root formula.
All counts are over labeled exponent assignments.
\end{abstract}

\noindent\textbf{Keywords:} constacyclic codes; repeated-root codes;
Galois hull; joint enumerator; transfer matrix; affine algebra.

\noindent\textbf{Mathematics Subject Classification:} 94B15, 94B05, 11T71.

% =====================================================================
\section{Introduction}
% =====================================================================

For a linear code \(C\) over a finite field, the hull --- the intersection
of \(C\) with its dual --- controls several constructions, including the
amount of entanglement needed in entanglement-assisted quantum codes and
the behaviour of LCD (linear complementary dual) codes. Replacing the
Euclidean inner product by a Galois inner product~\cite{FZ2017} leads to the
\(k\)-Galois hull \(\Hull_k(C) = C \cap D(C)\), where \(D(C)\) is the dual
with respect to a Frobenius-twisted pairing. Galois hulls of constacyclic
codes have been determined, at the level of hull dimensions and hull
generators, in a number of settings. Over finite fields,
Debnath, Prakash and Islam determined the Galois hull dimension of
constacyclic codes and the number of such codes with a given Galois hull
dimension \cite{DPI2023} (see also the correction \cite{DPI2023corr}); average Galois hull dimensions over finite
fields and over an auxiliary ring were computed by Debnath and Prakash
\cite{DP2025}, and constacyclic codes with small Galois hull dimension
were studied by Debnath, Islam, Yadav and Prakash \cite{DIYP2026}. Over
affine algebra rings, Debnath, Islam, Mart\'inez-Moro and Prakash obtained
Galois hull formulas for constacyclic codes \cite{DIMP2026}, working
mainly in a simple-root (idempotent) framework with a repeated-root case
under additional hypotheses on the order of the twist. Eldin and Sol\'e
developed generator polynomial matrices for the Galois hulls of
multi-twisted codes, of which constacyclic codes are the index-one case
\cite{ES2026}. Galois hulls of cyclic serial codes over a finite chain ring are studied in
\cite{Talbi2021}, and Galois hulls have likewise been used for quantum-code
constructions over a finite non-chain ring \cite{ZKZ2026}; repeated-root constacyclic codes and their (Euclidean)
duals, including LCD and self-dual enumerations in special lengths, have
a substantial literature \cite{CDL2015, Sharma2015, WZLL2023, ZWLL2025}.

The present paper studies a distributional question in the repeated-root
setting. Fix a square-free affine algebra
\(A \cong \prod_{s=1}^{N} K_s\) and a length \(n = P n_0\) in which the
characteristic divides \(n\) through the prime power \(P = p^{\varrho}\),
while \(\gcd(n_0, p) = 1\). Every constacyclic code over \(A\) is then
described, component by component, by exponents
\(a_f \in \{0, \dots, P\}\) attached to the distinct irreducible factors
of a separable polynomial \(x^{n_0} - \mu_s\). Rather than the dimension of
a single hull, we determine the \emph{joint distribution} of the pair
\((\dim_{\Fq} C, \dim_{\Fq} \Hull_k(C))\) over all labeled exponent
assignments, for compatible twists. The joint law is captured by a
transfer-matrix enumerator: each orbit of the reciprocal involution on
the factor set contributes the trace of a power of a \((P{+}1)\)-state
matrix whose entries record, edge by edge, the code-dimension and
hull-dimension contributions (Theorem~\ref{thm:T6}). From the enumerator
we extract the total number of codes, the code-dimension distribution,
the number of LCD codes, characterizations and counts of self-orthogonal
and dual-containing codes, and exact first and second moments of the hull
dimension. A point that requires care is the case of an orbit of length
one: its hull edge is a self-loop \(\min(P-A, A)\), whose mean
\(\nu_P = \lfloor P^2/4 \rfloor / (P+1)\) differs from the mean
\(\mu_P = P(2P+1)/(6(P+1))\) of an ordinary edge; conflating the two
gives an incorrect hull mean as soon as a fixed factor of the reciprocal
involution is present.

The paper is organized as follows. Section~\ref{sec:prelim} fixes the
conventions, in particular the code-first Galois pairing on which the
dual formulas depend. Sections~\ref{sec:param}--\ref{sec:hull} develop
the parametrization, the Galois dual, and the hull formula in exponent
form. Section~\ref{sec:enumerator} proves the transfer-matrix enumerator
and its specialization at \(P=1\). Sections~\ref{sec:corollaries} and
\ref{sec:moments} contain the counting corollaries and the moment
formulas. Section~\ref{sec:examples} gives worked examples and describes
the computational verification; Section~\ref{sec:limitations} records
the scope and limitations. Throughout, counts are over labeled exponent
assignments, not over equivalence or isometry classes of codes.

% =====================================================================
\section{Preliminaries and conventions}\label{sec:prelim}
% =====================================================================

Let \(p\) be a prime, \(q = p^e\), and let \(0 \le k < e\) be a fixed
Galois index. Let
\[
A \;\cong\; \prod_{s=1}^{N} K_s, \qquad
K_s = \mathbb{F}_{q^{m_s}} = \mathbb{F}_{p^{d_s}}, \quad d_s = e m_s,
\]
with fixed component labels. (Standard terminology for constacyclic codes
follows \cite{HP2003, MS1977}.) On \(K_s\) write
\(\sigma(\alpha) = \alpha^{p^k}\) and
\(\rho(\alpha) = \alpha^{p^{d_s - k}} = \sigma^{-1}\).

\begin{definition}[code-first Galois pairing]
For one component (the subscript \(s\) suppressed), the \(k\)-Galois
pairing on \(K^n\) is
\(\langle x, y \rangle_k = \sum_{i=0}^{n-1} x_i \, \sigma(y_i)\),
the \emph{code-first} dual of a code \(C \subseteq K^n\) is
\(D(C) = \{ y \in K^n : \langle c, y \rangle_k = 0 \text{ for all } c \in C\}\),
and the \(k\)-Galois hull is \(\Hull_k(C) = C \cap D(C)\).
The Euclidean dual is denoted \(C^\perp\).
\end{definition}

This is the \(k\)-Galois inner product introduced by Fan and
Zhang~\cite{FZ2017}. The slot order is part of the definition: the candidate \(y\) occupies
the \(\sigma\)-twisted slot, and this is what makes
\(D(C) = \rho(C^\perp)\) (coordinatewise \(\rho\);
Proposition~\ref{prop:transport}) rather than the analogous statement
with \(\sigma\). All dual formulas below use this convention. The global
pairing on \(A^n\) is the tuple of component pairings; a code over \(A\)
is an \(A\)-submodule of \(A^n\).

Fix \(\varrho \ge 0\), \(P = p^{\varrho}\), and \(n = P n_0\) with
\(\gcd(n_0, p) = 1\). Let \(\lambda = (\lambda_s) \in A^*\) and
\(M_s(x) = x^n - \lambda_s \in K_s[x]\). Since \(a \mapsto a^P\) is a
bijection of the finite field \(K_s\), there is a unique \(\mu_s \in K_s^*\)
with \(\mu_s^P = \lambda_s\). Writing \(\bar M_s(x) = x^{n_0} - \mu_s\), we
have \(M_s = \bar M_s^{\,P}\) in characteristic \(p\), and \(\bar M_s\) is
separable because \(\bar M_s' = n_0 x^{n_0-1}\) and every root of
\(\bar M_s\) is nonzero. Let \(\mathcal{F}_s\) be the set of distinct
monic irreducible factors of \(\bar M_s\) over \(K_s\); then
\(M_s = \prod_{f \in \mathcal{F}_s} f^P\).

An \emph{exponent assignment} on component \(s\) is a tuple
\(a = (a_f)_{f \in \mathcal{F}_s}\) with \(a_f \in \{0, \dots, P\}\); put
\(g_a = \prod_f f^{a_f}\) and \(C_s(a) = \langle g_a \rangle \subseteq
K_s[x]/(M_s)\). A global assignment is a tuple of component assignments.

A \emph{\(\lambda\)-constacyclic code} over \(A\) is an \(A\)-submodule of
\(A^n\) closed under the componentwise \(\lambda\)-constacyclic shift
\((c_0,\dots,c_{n-1}) \mapsto (\lambda c_{n-1}, c_0, \dots, c_{n-2})\).
Under the coefficient identification \(A^n \cong A[x]/(x^n-\lambda)\), this
shift is multiplication by \(x\); hence the \(\lambda\)-constacyclic codes
over \(A\) are exactly the ideals of \(A[x]/(x^n-\lambda)\), an equivalence
that needs no separability or compatibility hypothesis. From now on,
``code over \(A\)'' means a \(\lambda\)-constacyclic code over \(A\);
componentwise these are tuples of \(\lambda_s\)-constacyclic codes
(Proposition~\ref{prop:assembly}).

For a monic polynomial \(f(x) = \sum_{i=0}^{r} c_i x^i\) with
\(f(0) = c_0 \ne 0\), define the (Frobenius-twisted) reciprocal
\[
f^{\#}(x) \;=\; \rho(c_0)^{-1} \sum_{i=0}^{r} \rho(c_i)\, x^{\,r-i},
\]
which is monic of degree \(r\). Every \(f \in \mathcal{F}_s\) satisfies
\(f(0) \ne 0\) since \(\bar M_s(0) = -\mu_s \ne 0\).

\begin{definition}[compatibility and the involution \(\tau_s\)]
The twist \(\lambda_s\) is \emph{compatible} if
\(\lambda_s^{\,1 + p^{d_s - k}} = 1\).
Under compatibility, \(\tau_s : \mathcal{F}_s \to \mathcal{F}_s\),
\(\tau_s(f) = f^{\#}\), is a permutation (Theorem~\ref{thm:T2}). We
assume compatibility for every \(s\) from Section~\ref{sec:dual} on.
A \(\tau_s\)-orbit is written \(O = (f_0, \dots, f_{L-1})\) with
\(\tau_s(f_i) = f_{i+1}\), indices modulo \(L = L_O\); we set
\(\alpha_i = a_{f_i}\) and call \(\alpha_{i-1}\) the predecessor of
\(\alpha_i\). Since \(\#\) preserves degrees, all factors in an orbit
have the same degree; the \emph{weight} of \(O\) and of a factor \(f\) is
\(w_O = m_s \deg_{K_s} f\) and \(w_f = m_s \deg_{K_s} f\).
\end{definition}

Dimensions over \(\Fq\) are recovered from dimensions over \(K_s\) by
\(\dim_{\Fq} V = m_s \dim_{K_s} V\). We write \(u, z\) for indeterminates
and work in \(\Z[u,z]\).

% =====================================================================
\section{Parametrization and the reciprocal operation}\label{sec:param}
% =====================================================================

\begin{theorem}[T1: parametrization]\label{thm:T1}
The map \(a \mapsto C_s(a) = \langle g_a \rangle\) is a bijection from
\(\{0,\dots,P\}^{\mathcal{F}_s}\) onto the set of \(\lambda_s\)-constacyclic
codes of length \(n\) over \(K_s\), and
\(\dim_{K_s} C_s(a) = n - \sum_{f} a_f \deg f\).
Consequently there are \((P+1)^{\sum_s |\mathcal{F}_s|}\) labeled global
assignments, and \(\lambda\)-constacyclic codes over \(A\) correspond to tuples
of component assignments (Proposition~\ref{prop:assembly}).
\end{theorem}

\begin{proof}
Ideals of \(K_s[x]/(M_s)\) correspond to ideals \((M_s) \subseteq I\) of
the principal ideal domain \(K_s[x]\), hence to monic divisors
\(\bar g\) of \(M_s\); by unique factorization and \(M_s =
\prod f^P\), monic divisors are exactly the \(g_a\). For a monic divisor
\(g\) of degree \(r\), multiplication by \(g\) identifies
\(\langle g \rangle \subseteq K_s[x]/(M_s)\) with
\(K_s[x]/((M_s)/g)\), of dimension \(n - r\). No separability hypothesis
is used.
\end{proof}

\begin{theorem}[T2: the reciprocal operation]\label{thm:T2}
\begin{enumerate}
\item[(i)] \(\#\) is multiplicative on monic polynomials with nonzero
constant term; hence
\(\big(\prod_f f^{a_f}\big)^{\#} = \prod_f (f^{\#})^{a_f}\) for all
exponents \(a_f \ge 0\).
\item[(ii)] If \(f(\alpha) = 0\) with \(\alpha \ne 0\) in an algebraic
closure, then \(f^{\#}\big(\alpha^{-p^{d_s-k}}\big) = 0\); in fact
\(f^{\#}\) is the minimal polynomial of \(\alpha^{-p^{d_s-k}}\).
\item[(iii)] \(\lambda_s^{1+p^{d_s-k}} = 1 \iff
\mu_s^{1+p^{d_s-k}} = 1\). Under compatibility, \(\tau_s\) permutes
\(\mathcal{F}_s\) and preserves degrees.
\end{enumerate}
\end{theorem}

\begin{proof}
Let \(\bar\rho\) be the coefficientwise extension of \(\rho\) to
\(K_s[x]\); it is a ring automorphism with inverse \(\bar\sigma\).
For (i), \(\overline\rho(fg) = \bar\rho(f)\bar\rho(g)\), reversal
\(x^{r}f(1/x)\) is multiplicative in fixed degrees, and the normalizing
constants multiply since \((fg)_0 = f_0 g_0\); combining the three gives
\((fg)^{\#} = f^{\#} g^{\#}\), and induction on the number of factors
gives the displayed identity. For (ii), apply the Frobenius map
\(t \mapsto t^{p^{d_s-k}}\) of the algebraic closure to
\(\sum c_i \alpha^i = 0\) to get \(\sum \rho(c_i) \beta^{-i} = 0\) for
\(\beta = \alpha^{-p^{d_s-k}}\); multiplying by \(\beta^r\) yields
\(\rho(c_0) f^{\#}(\beta) = 0\). The same Frobenius transport preserves
degrees of minimal polynomials, so \(f^{\#}\) is the (monic irreducible)
minimal polynomial of \(\beta\). For (iii),
\((\mu_s^{1+p^{d_s-k}})^P = \lambda_s^{1+p^{d_s-k}}\) and in
characteristic \(p\), \(y^P = 1\) implies \(y = 1\); now (ii) shows
\(\tau_s(f)\) is the minimal polynomial of a root \(\beta\) with
\(\beta^{n_0} = \rho(\mu_s)^{-1} = \mu_s\) (using compatibility in the
form \(\rho(\mu_s) = \mu_s^{-1}\)), so \(\tau_s(f) \mid \bar M_s\);
bijectivity follows by applying the same construction with \(\sigma\) in
place of \(\rho\), which gives the inverse map.
\end{proof}

% =====================================================================
\section{The Galois dual}\label{sec:dual}
% =====================================================================

The first step is a Euclidean-dual statement that uses no separability
at all.

\begin{lemma}[Euclidean dual generator]\label{lem:RR3}
Let \(g\) be a monic divisor of \(x^n - \lambda\), \(\lambda \in K^*\),
\(h = (x^n-\lambda)/g\), \(r = \deg g\), \(s = n - r\), and
\(C = \langle g \rangle\) under the coefficient identification. Then
\(h_0 \ne 0\), and with
\(h^{\natural}(x) = h_0^{-1} \sum_{i=0}^{s} h_i x^{s-i}\) (the ordinary
reciprocal, without Frobenius),
\(C^\perp = \langle h^{\natural} \rangle\), a \(\lambda^{-1}\)-constacyclic
code.
\end{lemma}

\begin{proof}
From \(gh = x^n - \lambda\), \(g_0 h_0 = -\lambda \ne 0\). Both
\(\langle h^{\natural} \rangle\) and \(C^\perp\) have dimension \(r\):
the first by the divisor dimension count of Theorem~\ref{thm:T1}
(applied to any modulus), the second by rank--nullity. The boundary
cases are immediate from the definitions: if \(r = 0\), then \(g = 1\),
\(C = K^n\) and \(C^\perp = \{0\}\), while \(h^{\natural}\) has degree
\(n\) and hence no nonzero multiple of degree \(< n\); if \(s = 0\), then
\(h = 1\), \(C = \{0\}\) and \(C^\perp = K^n\), while
\(h^{\natural} = 1\). Assume \(r, s \ge 1\) below, so that \(V\) is
defined (and \(v = 0\) gives \(d = 0\)). We show orthogonality. For polynomials \(a, d\) of degree \(< n\), write
\(\operatorname{rev}_n(d) = x^{n-1} d(1/x)\); then
\(\sum_i a_i d_i = [x^{n-1}]\big(a(x) \operatorname{rev}_n(d)(x)\big)\).
Take \(a = gu\), \(\deg u < s\), and \(d = h^{\natural} v\),
\(\deg v < r\). Since
\(h^{\natural}(1/x) = h_0^{-1} x^{-s} h(x)\), one computes
\(\operatorname{rev}_n(d) = h_0^{-1} h(x) V(x)\) with
\(V(x) = x^{r-1} v(1/x)\), \(\deg V \le r-1\). Hence
\(a(x)\operatorname{rev}_n(d)(x) = h_0^{-1} (x^n - \lambda) u(x) V(x)\)
with \(\deg(uV) \le n-2\); the coefficient of \(x^{n-1}\) is \(0\).
Thus \(\langle h^{\natural}\rangle \subseteq C^\perp\), and equality
follows by dimensions. That \(h^{\natural} \mid x^n - \lambda^{-1}\)
follows from \((gh)^* = g^* h^* = -\lambda(x^n - \lambda^{-1})\) for the
unnormalized reciprocal \(\phi^*(x) = x^{\deg\phi}\phi(1/x)\).
\end{proof}

\begin{proposition}[transport by \(\rho\)]\label{prop:transport}
Let \(K = \mathbb{F}_{p^d}\) with \(\sigma, \rho\) as in
Section~\ref{sec:prelim}. For every \(K\)-linear code \(C \subseteq K^n\),
\(D(C) = \rho(C^\perp)\) (coordinatewise). Consequently, for
\(C = \langle g \rangle \subseteq K[x]/(x^n-\lambda)\) and
\(h = (x^n-\lambda)/g\),
\(D(C) = \langle h^{\#} \rangle \subseteq
K[x]/(x^n - \lambda^{-p^{d-k}})\); if \(\lambda\) is compatible, the
modulus is again \(x^n - \lambda\).
\end{proposition}

\begin{proof}
By the code-first definition, \(y \in D(C)\) iff
\(\sum_i c_i \sigma(y_i) = 0\) for all \(c \in C\), i.e.\ iff
\(\sigma(y) \in C^\perp\) (coordinatewise), i.e.\ iff
\(y \in \rho(C^\perp)\) since \(\rho = \sigma^{-1}\); this is exactly
where the slot order enters. Coefficientwise application of an
automorphism sends the ideal \(\langle e \rangle\) in
\(K[x]/(x^n-\eta)\) to \(\langle \bar\rho(e)\rangle\) in
\(K[x]/(x^n - \rho(\eta))\); applying this to Lemma~\ref{lem:RR3} and
comparing \(\bar\rho(h^{\natural})\) with \(h^{\#}\) gives the generator.
Compatibility \(\lambda^{1+p^{d-k}}=1\) gives
\(\lambda^{-p^{d-k}} = \lambda\).
\end{proof}

\begin{theorem}[T3: dual exponents]\label{thm:T3}
For every assignment \(a\), with \(h_a = M_s / g_a\),
\(D(C_s(a)) = \langle h_a^{\#} \rangle\); under compatibility this is a
\(\lambda_s\)-constacyclic code in the same quotient. On a
\(\tau_s\)-orbit \(O = (f_0,\dots,f_{L-1})\) with \(\alpha_i = a_{f_i}\),
the exponent of \(f_i\) in \(h_a^{\#}\) is
\[
\beta_i \;=\; P - \alpha_{i-1} \qquad (\text{indices mod } L).
\]
\end{theorem}

\begin{proof}
The generator statement is Proposition~\ref{prop:transport}. For the
exponents, \(h_a = \prod_{f} f^{P - a_f}\), so by Theorem~\ref{thm:T2}(i)
and the permutation property,
\(h_a^{\#} = \prod_f (\tau_s f)^{P-a_f} = \prod_{g} g^{P - a_{\tau_s^{-1} g}}\);
for \(g = f_i\), \(\tau_s^{-1}(f_i) = f_{i-1}\), giving \(P - \alpha_{i-1}\).
\end{proof}

\begin{remark}
Without compatibility, \(D(C_s(a)) = \langle h_a^{\#}\rangle\) and the
exponent formula remain valid, but in the modulus
\(x^n - \lambda_s^{-p^{d_s-k}}\). Compatibility is what allows the dual
to be viewed inside the same constacyclic family; it cannot simply be
dropped (e.g.\ in \(\mathbb{F}_9\) with \(k=1\), a twist \(\lambda\) of
order \(8\) satisfies \(\lambda^{1+3} = \lambda^4 = -1 \ne 1\) and indeed
\(\lambda^{-3} \ne \lambda\)).
\end{remark}

% =====================================================================
\section{The Galois hull}\label{sec:hull}
% =====================================================================

\begin{lemma}[intersection of ideals]\label{lem:lcm}
For monic divisors \(g, l\) of a monic \(M \in K[x]\), the images of
\((g)\) and \((l)\) in \(K[x]/(M)\) intersect in the image of
\((\lcm(g,l))\), and
\(\lcm(g,l) = \prod_{\pi} \pi^{\max(v_\pi(g), v_\pi(l))}\).
\end{lemma}

\begin{proof}
In \(K[x]\), \((g) \cap (l) = (\lcm(g,l))\) since both consist of the
common multiples of \(g\) and \(l\), characterized valuation-wise by the
displayed maximum. If \(y\) lies in both images in the quotient, any
lift \(Y\) satisfies \(Y \in (g)\) and \(Y \in (l)\) because
\((M) \subseteq (g), (l)\); hence \(Y \in (\lcm(g,l))\). The reverse
inclusion is immediate from \(g, l \mid \lcm(g,l)\).
\end{proof}

\begin{theorem}[T4: hull generator and dimensions]\label{thm:T4}
Under compatibility, for every assignment \(a\),
\(\Hull_k(C_s(a)) = \langle \lcm(g_a, h_a^{\#}) \rangle\), the exponent
of \(f_i\) in the hull generator being \(\max(\alpha_i, P-\alpha_{i-1})\).
Moreover
\[
\dim_{\Fq} C(a) \;=\; \sum_{s} \sum_{O} w_O \sum_{i} (P - \alpha_i),
\qquad
\dim_{\Fq} \Hull_k(C(a)) \;=\; \sum_{s} \sum_{O} w_O \sum_i
\min(P-\alpha_i,\, \alpha_{i-1}),
\]
the inner sums over orbit positions (indices mod \(L_O\)); here and
below, \(\alpha_i\) refers to the orbit \((s,O)\) under discussion.
\end{theorem}

\begin{proof}
Work in one component. By Theorems~\ref{thm:T1} and \ref{thm:T3},
\(C(a)\) and \(D(C(a))\) are the images of \((g_a)\) and
\((h_a^{\#})\), both dividing \(M = x^n-\lambda\); Lemma~\ref{lem:lcm}
gives the hull generator, whose \(f_i\)-exponent is
\(\max(\alpha_i, \beta_i) = \max(\alpha_i, P-\alpha_{i-1})\). Dimensions
are \(n - \deg(\text{generator})\) over \(K\); grouping factors by
orbits (constant degree on an orbit) and using
\(P - \max(u,v) = \min(P-u, P-v)\) with \(v = P-\alpha_{i-1}\) gives the
hull formula; multiplication by \(m_s\) and summation over \(s\) gives
the \(\Fq\)-dimensions.
\end{proof}

% =====================================================================
\section{The transfer-matrix enumerator}\label{sec:enumerator}
% =====================================================================

For \(P \ge 1\), \(w \ge 1\), define the \((P{+}1)\times(P{+}1)\) matrix
over \(\Z[u,z]\)
\[
T^{(P)}_w(u,z)_{a,b} \;=\; u^{\,w(P-b)}\, z^{\,w \min(P-b,\, a)},
\qquad a, b \in \{0,\dots,P\},
\]
the row index \(a\) being the previous exponent and the column index
\(b\) the current one.

\begin{theorem}[T5: local transfer trace]\label{thm:T5}
For a \(\tau_s\)-orbit \(O\) of length \(L\) and weight \(w\),
\[
\tr\!\big(T^{(P)}_w(u,z)^L\big)
\;=\; \sum_{\alpha \in \{0,\dots,P\}^L}
u^{\,w \sum_i (P-\alpha_i)}\,
z^{\,w \sum_i \min(P-\alpha_i,\, \alpha_{i-1})},
\quad \text{indices mod } L.
\]
\end{theorem}

\begin{proof}
Expanding the matrix power,
\((T^L)_{a_0,a_L} = \sum_{a_1,\dots,a_{L-1}} \prod_{j=1}^{L} T_{a_{j-1},a_j}\);
setting \(a_L = a_0\) and summing over \(a_0\) gives
\(\tr(T^L) = \sum_{a_0,\dots,a_{L-1}} \prod_{j=1}^{L} T_{a_{j-1},a_j}\).
These closed walks are in bijection with assignments
\(\alpha = (\alpha_0,\dots,\alpha_{L-1})\) (orbit positions are labeled;
a walk and its rotation are different assignments). The \(j\)-th edge
contributes \(u^{w(P-a_j)} z^{w\min(P-a_j, a_{j-1})}\); summing exponents
over the walk and reindexing gives the claim. For \(L=1\) the single edge
is the self-loop \(a_0 \to a_0\) and the formula reads
\(\sum_a u^{w(P-a)} z^{w\min(P-a,a)}\), as required.
\end{proof}

\begin{proposition}[product assembly]\label{prop:assembly}
Ideals of \(A[x]/(x^n-\lambda) \cong \prod_s K_s[x]/(x^n-\lambda_s)\) are
exactly the products of component ideals; the global dual and hull
decompose componentwise; and
\(\dim_{\Fq} C = \sum_s m_s \dim_{K_s} C_s\), likewise for the hull.
\end{proposition}

\begin{proof}
The ring isomorphism is coefficientwise; quotienting a product by the
product ideal \((x^n-\lambda_s)_s\) gives the product of quotients (this
is the direct-product structure of the square-free algebra \(A\) itself;
no coprimality between components is involved). Ideals of a finite
product of rings are products of ideals (project, and recover the
components with the idempotents of \(A\)). The pairing is componentwise
and vanishes exactly when every component does, giving the dual and hull
statements; dimensions add over the direct sum \(A^n \cong \bigoplus_s
K_s^n\) with the factor \([K_s:\Fq] = m_s\).
\end{proof}

\begin{theorem}[T6: joint enumerator]\label{thm:T6}
Under the standing hypotheses (common \(P, n, k\); compatible twists;
labeled components),
\[
\mathcal{E}(u,z) \;:=\;
\sum_{\text{labeled } a} u^{\dim_{\Fq} C(a)} z^{\dim_{\Fq} \Hull_k(C(a))}
\;=\; \prod_{s=1}^{N} \prod_{O \in \mathcal{O}_s}
\tr\!\Big(T^{(P)}_{w_O}(u,z)^{L_O}\Big),
\]
where \(\mathcal{O}_s\) is the set of \(\tau_s\)-orbits.
The coefficient of \(u^K z^H\) is the number of labeled assignments with
\(\dim_{\Fq} C = K\) and \(\dim_{\Fq} \Hull_k(C) = H\).
\end{theorem}

\begin{proof}
By Theorem~\ref{thm:T1} and Proposition~\ref{prop:assembly}, a global
assignment is the disjoint union of its orbit restrictions
\(\alpha^{(O)} \in \{0,\dots,P\}^{L_O}\), and by Theorem~\ref{thm:T4}
each of the two dimensions is a sum over \((s,O)\) of a quantity
depending only on \(\alpha^{(O)}\). The finite distributive law
(expanding a product of sums over a Cartesian product index set)
factorizes the global sum into the product of the orbit sums
\(\mathcal{E}_O(u,z) = \sum_{\alpha} u^{w_O \sum (P-\alpha_i)}
z^{w_O \sum \min(P-\alpha_i,\alpha_{i-1})}\), and Theorem~\ref{thm:T5}
identifies \(\mathcal{E}_O(u,z) = \tr(T^{(P)}_{w_O}(u,z)^{L_O})\).
\end{proof}

\begin{remark}[hypotheses of the identity]
As an identity in \(\Z[u,z]\), Theorem~\ref{thm:T5} holds for every
integer \(P \ge 1\); the prime-power form \(P = p^{\varrho}\) is needed
for the field construction behind Theorems~\ref{thm:T1}--\ref{thm:T4},
not for the matrix identity. The enumerator counts labeled assignments;
the trace formula is insensitive to reversing the orbit orientation
(the distribution is unchanged under \(\alpha_i \mapsto \alpha_{-i}\)),
so the predecessor convention \(\alpha_{i-1}\) in
Theorem~\ref{thm:T4} is an essential part of the statement, fixed by the
code-first pairing of Section~\ref{sec:prelim}.
\end{remark}

\begin{proposition}[simple-root recovery at \(P=1\)]\label{prop:P1}
At \(\varrho = 0\) (\(P=1\)), assignments are binary, \(g_a\) is a
squarefree divisor, the hull edge term \(\min(1-\alpha_i, \alpha_{i-1})\)
is the indicator of the transition \(1 \to 0\) along the orbit, and
\[
T^{(1)}_w(u,z) \;=\;
\begin{pmatrix} u^w & 1 \\ u^w z^w & 1 \end{pmatrix}
\qquad (\text{rows/columns } a,b \in \{0,1\}),
\]
so Theorem~\ref{thm:T6} specializes to the binary-state simple-root
joint enumerator with the same conventions.
\end{proposition}

\begin{proof}
For \(P=1\), \(\min(1-b, a)\) is \(1\) exactly for \((a,b) = (1,0)\);
reading off the four entries gives the matrix, and \(\mu_s = \lambda_s\),
\(\bar M_s = M_s\) gives the squarefree parametrization.
\end{proof}

% =====================================================================
\section{Counting corollaries}\label{sec:corollaries}
% =====================================================================

\begin{corollary}[C1: total count]\label{cor:C1}
\(\mathcal{E}(1,1) = (P+1)^{\sum_s |\mathcal{F}_s|}\).
\end{corollary}

\begin{proof}
At \(u=z=1\) each \(T^{(P)}_w\) is the all-ones matrix \(J\) with
\(J^L = (P+1)^{L-1} J\), so \(\tr(J^{L_O}) = (P+1)^{L_O}\); product over
orbits gives \((P+1)^{\sum_s \sum_O L_O} = (P+1)^{\sum_s |\mathcal{F}_s|}\).
\end{proof}

\begin{corollary}[C2: code-dimension distribution]\label{cor:C2}
\(\displaystyle
\mathcal{E}(u,1) = \prod_{s} \prod_{f \in \mathcal{F}_s}
\big(1 + u^{w_f} + \dots + u^{P w_f}\big)\).
\end{corollary}

\begin{proof}
At \(z=1\) the entry \(T_{a,b} = u^{w(P-b)}\) is independent of the row,
so \(T = c\,r\) is rank one with \(r \cdot c = \sum_{j=0}^{P} u^{wj}\);
then \(\tr(T^L) = (r\cdot c)^{L-1} \tr(T) = \big(\sum_{j=0}^P u^{wj}\big)^L\).
Within an orbit all \(w_f\) equal \(w_O\), giving
\(\prod_{f \in O} (1 + u^{w_f} + \dots + u^{Pw_f})\) per orbit.
\end{proof}

\begin{corollary}[C3: LCD count]\label{cor:C3}
The number of labeled assignments with \(\Hull_k(C(a)) = 0\) is
\(2^{\#\text{orbits}}\), independent of \(P\).
\end{corollary}

\begin{proof}
The hull dimension is a sum of nonnegative terms
\(w_O \min(P-\alpha_i, \alpha_{i-1})\); it vanishes iff every orbit
satisfies \(\min(P-\alpha_i, \alpha_{i-1}) = 0\) for all \(i\). On one
orbit, if some \(\alpha_m \notin \{0,P\}\) (so \(P \ge 2\)), then
\(\min(P-\alpha_{m+1}, \alpha_m) = 0\) with \(\alpha_m > 0\) forces
\(\alpha_{m+1} = P\), and inductively every \(\alpha_i = P\) around the
cycle, a contradiction (for \(L=1\), \(\min(P-\alpha_0,\alpha_0) = 0\)
already forces \(\alpha_0 \in \{0,P\}\)). Hence each \(\alpha_i \in
\{0,P\}\); if both values occurred, at a boundary position with
\(\alpha_i = 0\), \(\alpha_{i-1} = P\) the term would be
\(\min(P,P) = P \ge 1\). So exactly two tuples per orbit
(\(\alpha \equiv 0\) or \(\alpha \equiv P\)), independently across
orbits.
\end{proof}

\begin{corollary}[C4: self-orthogonal and dual-containing codes]
\label{cor:C4}
\(C(a) \subseteq D(C(a))\) if and only if
\(\alpha_{i-1} + \alpha_i \ge P\) on every orbit edge, and then the count
is \(\prod_O \tr(B_P^{L_O})\) with \((B_P)_{a,b} = 1\) iff \(a+b \ge P\).
Dually, \(D(C(a)) \subseteq C(a)\) if and only if
\(\alpha_{i-1} + \alpha_i \le P\) on every edge, with count
\(\prod_O \tr((B'_P)^{L_O})\), \((B'_P)_{a,b} = 1\) iff \(a+b \le P\).
The two counts are equal.
\end{corollary}

\begin{proof}
Since \(\Hull = C \cap D(C) \subseteq C\), \(C \subseteq D(C)\) iff
\(\dim \Hull = \dim C\). By Theorem~\ref{thm:T4},
\(\dim C - \dim \Hull = \sum_{s,O} w_O \sum_i \max(0, P-\alpha_i-\alpha_{i-1})\)
because \((P-\alpha_i) - \min(P-\alpha_i,\alpha_{i-1}) =
\max(0, P-\alpha_i-\alpha_{i-1})\); every summand is nonnegative, so the
total vanishes iff every edge satisfies \(\alpha_{i-1}+\alpha_i \ge P\).
The admissible tuples on one orbit are the closed walks of the \(0\)-\(1\)
matrix \(B_P\), counted by \(\tr(B_P^{L_O})\) as in
Theorem~\ref{thm:T5}. For the dual side, \(\dim D(C) =
\sum w_O \sum_i \alpha_{i-1} = \sum w_O \sum_i \alpha_i\) by
Theorem~\ref{thm:T3}, and \(\alpha_{i-1} - \min(P-\alpha_i,\alpha_{i-1})
= \max(0, \alpha_i+\alpha_{i-1}-P)\); the same argument gives the
condition \(\alpha_{i-1}+\alpha_i \le P\) and the count
\(\tr((B'_P)^{L_O})\). The involution \(\alpha_i \mapsto P-\alpha_i\)
maps one edge condition to the other, so the counts agree.
\end{proof}

% =====================================================================
\section{Moments of the hull dimension}\label{sec:moments}
% =====================================================================

Choose a labeled assignment uniformly at random; equivalently, all
exponents \(A_i\) (over all components, orbits, positions) are
independent and uniform on \(\{0,\dots,P\}\). For one orbit write
\(X_i = \min(P-A_i, A_{i-1})\) (indices mod \(L\)) and
\(S_L = \sum_i X_i\); the orbit contributes \(w_O S_{L_O}\) to
\(\dim_{\Fq}\Hull_k\) (Theorem~\ref{thm:T4}). Set
\[
\mu_P = \frac{P(2P+1)}{6(P+1)},
\qquad
\nu_P = \frac{\lfloor P^2/4 \rfloor}{P+1}.
\]

\begin{lemma}[edge moments]\label{lem:edge}
For independent \(U, V\) uniform on \(\{0,\dots,P\}\),
\(\E[\min(U,V)] = \mu_P\) and
\(\E[\min(U,V)^2] = \frac{P(P^2+P+1)}{6(P+1)}\).
For \(0 \le c \le P\), \(F(c) := \E[\min(c,U)] =
\frac{c(2P-c+1)}{2(P+1)}\).
\end{lemma}

\begin{proof}
For a nonnegative integer variable,
\(\E[X] = \sum_{t\ge1} \Pr(X \ge t)\) and
\(\E[X^2] = \sum_{t\ge1} (2t-1)\Pr(X\ge t)\).
Here \(\Pr(\min(U,V) \ge t) = ((P-t+1)/(P+1))^2\) for \(1 \le t \le P\);
summing via \(\sum_{r=1}^P r^2 = P(P+1)(2P+1)/6\) and
\(\sum r^3 = (P(P+1)/2)^2\) gives both formulas. The last claim is
\(F(c) = \sum_{t=1}^c (P-t+1)/(P+1)\).
\end{proof}

\begin{theorem}[C5: means]\label{thm:C5}
\(\displaystyle \E[\dim_{\Fq} C] = \frac{n}{2} \sum_s m_s\), and
\[
\E[\dim_{\Fq} \Hull_k(C)] \;=\;
\nu_P \sum_{O:\, L_O = 1} w_O \;+\;
\mu_P \sum_{O:\, L_O \ge 2} w_O L_O .
\]
\end{theorem}

\begin{proof}
\(\E[P - A_f] = P/2\) and \(\sum_{f \in \mathcal{F}_s} \deg f = n_0\)
give the first formula via Theorem~\ref{thm:T4}. For the hull, linearity
over orbits reduces to \(\E[S_L]\). If \(L \ge 2\), each \(X_i\) is the
minimum of two independent uniforms (for \(L=2\), \(X_0, X_1\) are the
two ordered pairings of the same two distinct variables), so
\(\E[S_L] = L\mu_P\) by Lemma~\ref{lem:edge}. If \(L = 1\),
\(\E[S_1] = \frac{1}{P+1}\sum_{a=0}^P \min(P-a,a)\); the summands are
symmetric under \(a \mapsto P-a\), giving \(P^2/4\) for \(P = 2m\) and
\((P^2-1)/4\) for \(P = 2m+1\), i.e.\ \(\lfloor P^2/4 \rfloor\); hence
\(\E[S_1] = \nu_P\). Note \(\nu_1 = 0\): at \(P=1\) an \(L_O=1\) orbit
contributes nothing, and \(\E[\Hull] = \frac14 \sum_{L_O\ge2} w_O L_O\).
\end{proof}

For the variance, orbit contributions are independent (disjoint exponent
variables), so with \(V_L(P) := \operatorname{Var}(S_L)\),
\(\operatorname{Var}(\dim_{\Fq}\Hull_k) = \sum_{s,O} w_O^2 V_{L_O}(P)\);
note the squared weights, in contrast to the mean.

\begin{lemma}[variance pieces]\label{lem:varpieces}
\begin{enumerate}
\item[(i)] \(L=1\):
\(V_1(P) = \dfrac{P(P+2)(P^2+2P+4)}{48(P+1)^2}\) if \(P\) is even, and
\(V_1(P) = \dfrac{(P-1)(P+3)}{48}\) if \(P\) is odd.
\item[(ii)] For \(L \ge 3\), two adjacent edges satisfy
\(\operatorname{Cov}(X_i, X_{i+1}) = \kappa_P :=
-\dfrac{P(P+2)(7P^2+14P+9)}{360(P+1)^2} < 0\), while non-adjacent edge
pairs are independent.
\item[(iii)] For \(L=2\):
\(\E[X_0 X_1] = \dfrac{P^2(P-1)}{12(P+1)}\) and
\(\operatorname{Cov}(X_0,X_1) = -\dfrac{P^2(P+2)^2}{36(P+1)^2}\).
\end{enumerate}
\end{lemma}

\begin{proof}
(i) \(\E[\min(P-A,A)^2] = \frac{1}{P+1}\sum_a \min(P-a,a)^2\); for
\(P = 2m\) the squared values are \(0^2,\dots,(m-1)^2,m^2,(m-1)^2,\dots,0^2\),
summing to \(m(2m^2+1)/3\), and for \(P=2m+1\) each of \(0^2,\dots,m^2\)
occurs twice, summing to \(m(m+1)(2m+1)/3\); subtracting \(\nu_P^2\)
(Theorem~\ref{thm:C5}) and simplifying gives the claim.
(ii) \(X_i, X_{i+1}\) share exactly the variable \(A_i\). Conditioning on
\(A_i = b\) makes them conditionally independent with means
\(F(P-b)\) and \(F(b)\) (Lemma~\ref{lem:edge}), so
\(\E[X_iX_{i+1}] = \frac{1}{P+1}\sum_{b=0}^P F(P-b)F(b)
= \frac{P(P-1)(11P^2+15P+6)}{120(P+1)^2}\) by the power sums
\(\sum b^j\), \(j \le 4\); subtracting \(\mu_P^2\) gives \(\kappa_P\),
negative for \(P \ge 1\). Non-adjacent edges depend on disjoint pairs of
variables, hence are independent.
(iii) \(X_0 \ge s\) and \(X_1 \ge t\) hold simultaneously iff
\(s + t \le P\), in which case each of the two variables has
\(P-s-t+1\) admissible values; thus
\(\E[X_0X_1] = \sum_{s,t\ge1} \Pr(X_0\ge s, X_1\ge t)
= \frac{1}{(P+1)^2} \sum_{r=1}^{P-1} (P-r) r^2
= \frac{P^2(P-1)}{12(P+1)}\), and subtracting \(\mu_P^2\) gives the
covariance.
\end{proof}

\begin{theorem}[C6: variance of the hull dimension]\label{thm:C6}
\(\displaystyle
\operatorname{Var}\big(\dim_{\Fq} \Hull_k(C)\big)
= \sum_{s} \sum_{O \in \mathcal{O}_s} w_O^2\, V_{L_O}(P)\), where
\[
V_1(P) \text{ is as in Lemma~\ref{lem:varpieces}(i)}, \quad
V_2(P) = \frac{P(P+2)(P^2+2P+3)}{18(P+1)^2}, \quad
V_L(P) = \frac{L\, P(P+2)(P^2+2P+2)}{60(P+1)^2}\ \ (L \ge 3).
\]
\end{theorem}

\begin{proof}
Independence across orbits gives the weighted sum. For one orbit:
\(L=1\) is Lemma~\ref{lem:varpieces}(i). For \(L=2\),
\(\operatorname{Var}(S_2) = 2\operatorname{Var}(X) + 2\operatorname{Cov}(X_0,X_1)\);
with \(\operatorname{Var}(X) = \frac{P(P+2)(2P^2+4P+3)}{36(P+1)^2}\)
(from Lemma~\ref{lem:edge}) and (iii), the bracket factors as
\(P(P+2)(P^2+2P+3)\), giving \(V_2(P)\). For \(L \ge 3\),
\(\operatorname{Var}(S_L) = L\operatorname{Var}(X) + 2L\kappa_P\) (exactly
the \(L\) adjacent pairs contribute, by (ii)), and
\(\operatorname{Var}(X) + 2\kappa_P =
\frac{P(P+2)(P^2+2P+2)}{60(P+1)^2}\). At \(P=1\): \(V_1(1) = 0\),
\(V_2(1) = 1/4\), \(V_L(1) = L/16\). The \(L=2\) formula is not the
\(L \ge 3\) formula evaluated at \(L=2\) (at \(P=2\): \(44/81\) versus
\(8/27\)), because the two edges of a 2-orbit share both variables.
\end{proof}

% =====================================================================
\section{Examples and computational verification}\label{sec:examples}
% =====================================================================

\begin{example}[one orbit, \(P=2\), \(L=2\), \(w=1\), by hand]
\label{ex:E1}
The nine assignments \(\alpha = (\alpha_0,\alpha_1)\) give
\((\sum_i (2-\alpha_i), \sum_i \min(2-\alpha_i,\alpha_{i-1}))\) as
follows:
\((0,0)\mapsto(4,0)\), \((0,1)\mapsto(3,1)\), \((0,2)\mapsto(2,2)\),
\((1,0)\mapsto(3,1)\), \((1,1)\mapsto(2,2)\), \((1,2)\mapsto(1,1)\),
\((2,0)\mapsto(2,2)\), \((2,1)\mapsto(1,1)\), \((2,2)\mapsto(0,0)\).
Hence \(\E[S_2] = 10/9 = 2\mu_2\), \(\operatorname{Var}(S_2) = 44/81 =
V_2(2)\) (the \(L=2\) case of Theorem~\ref{thm:C6}), and exactly two
assignments have hull contribution \(0\), as Corollary~\ref{cor:C3}
requires per orbit.
\end{example}

\begin{example}[simple-root orbit, \(P=1\), \(L=3\), \(w=1\)]\label{ex:E2}
The joint distribution is
\(\{(0,0):1,\ (1,1):3,\ (2,1):3,\ (3,0):1\}\): the hull contribution
counts \(1\!\to\!0\) transitions (Proposition~\ref{prop:P1}), with
\(\E = 3/4 = 3\mu_1\) and variance \(3/16 = V_3(1)\).
\end{example}

\begin{example}[a repeated-root field case]\label{ex:E3}
Take \(K = \mathbb{F}_4\), \(k=1\), \(\lambda=1\), \(P=2\), \(n_0=5\),
\(n=10\). The factors of \(x^5-1\) over \(\mathbb{F}_4\) have degrees
\(1,2,2\), with \(\tau\)-orbit lengths \(1,2\); there are \(3^3=27\)
assignments. The joint \((\dim_K C, \dim_K \Hull_1(C))\) distribution,
computed directly from generator polynomials and the rank of the
code-first Gram matrix \(\big(\sum_j B_{t,j}\sigma(B_{r,j})\big)\), is
\[
\begin{aligned}
\{&(0,0):1, (1,1):1, (2,0):1, (2,2):2, (3,3):2, (4,2):2, (4,4):3,
(5,5):3, (6,2):2,\\
&(6,4):3, (7,3):2, (8,0):1, (8,2):2, (9,1):1, (10,0):1\},
\end{aligned}
\]
in agreement with Theorems~\ref{thm:T4} and \ref{thm:T6}. It gives
\(\E[\Hull] = 23/9 = \nu_2 + 2\cdot 2\mu_2 = \tfrac13 + \tfrac{20}{9}\)
(the length-one orbit contributes through \(\nu_2\), not \(\mu_2\)),
\(\operatorname{Var}(\Hull) = 194/81\) from Theorem~\ref{thm:C6}, LCD
count \(4 = 2^2\), and self-orthogonal \(=\) dual-containing count \(12\).
\end{example}

\begin{example}[orbit of length \(6\)]\label{ex:E4}
Take \(K = \mathbb{F}_8\), \(k=1\), \(\lambda=1\), \(P=2\), \(n_0=7\),
\(n=14\). The seven linear factors of \(x^7-1\) form \(\tau\)-orbits of
lengths \(1\) and \(6\); over the \(3^7 = 2187\) assignments the Gram-rank
computation again agrees with Theorems~\ref{thm:T4} and \ref{thm:T6}
assignment by assignment, with \(\E[\Hull] = 11/3\),
\(\operatorname{Var}(\Hull) = 10/9 = V_1(2) + V_6(2) = \tfrac29 +
\tfrac89\), LCD count \(4\), and self-orthogonal \(=\) dual-containing
count \(258\). Because the exponent formula is checked against the
actual hull dimensions of actual codes on an orbit of length \(6\), this
example also certifies the predecessor orientation \(\alpha_{i-1}\):
on such an orbit the predecessor and successor readings give different
hull dimensions for generic \(\alpha\), so the check would detect a
transposed convention (cf.\ the remark after Theorem~\ref{thm:T6}).
\end{example}

\begin{example}[two components]\label{ex:E5}
For \(A = \mathbb{F}_2 \times \mathbb{F}_4\) (\(q=2\), \(k=0\),
\(\lambda=(1,1)\), \(P=2\), \(n_0=7\)), the orbit data
\((L_O, w_O)\) are \((1,1), (2,3)\) and \((1,2), (2,6)\). Convolving the
per-orbit distributions of Theorem~\ref{thm:T5} gives the global joint
distribution over \(3^6 = 729\) labeled assignments, with LCD count
\(16 = 2^4\), self-orthogonal and dual-containing counts \(144\) each,
\(\E[\Hull] = 11\) (Theorem~\ref{thm:C5}), and
\(\operatorname{Var}(\Hull) = 230/9\) (Theorem~\ref{thm:C6}); the same
numbers result from adding independently computed component dimensions.
\end{example}

\paragraph{Computational verification.}
All formulas above were checked by exact (integer and rational)
computation in three independent ways: direct enumeration of exponent
tuples; the transfer traces of Theorem~\ref{thm:T5} evaluated both as
polynomial matrices and at integer points; and Gram-rank brute force
over explicitly constructed codes in the cases of
Examples~\ref{ex:E3}--\ref{ex:E5} and further cases
(\(\mathbb{F}_4\) with \(n_0=3\); \(\mathbb{F}_9\) with \(k=1\),
\(\lambda = \pm 1\), \(P=3\); \(\mathbb{F}_3\) with \(\lambda=-1\),
\(P=3\)), comprising \(2769\) field-level assignments in total, in
addition to an orbit grid of \(52\) pairs \((P,L)\),
\(P \in \{1,2,3,4,5,7,8,9\}\). The moment closed forms were also
re-derived symbolically from their defining sums. These computations are
supporting evidence for the tested cases; the proofs of
Sections~\ref{sec:param}--\ref{sec:moments} are independent of them.

% =====================================================================
\section{Limitations and scope}\label{sec:limitations}
% =====================================================================

\begin{itemize}
\item \textbf{Labeled counting.} All enumerations count labeled exponent
assignments over a fixed labeled decomposition
\(A \cong \prod_s K_s\). They are not counts of equivalence or isometry
classes of codes, and no such interpretation should be attached to
\(\mathcal{E}(u,z)\).
\item \textbf{Compatible twists only.} For incompatible \(\lambda_s\),
Theorem~\ref{thm:T3} persists in the two-modulus form
\(x^n - \lambda_s^{-p^{d_s-k}}\), but the hull formula and everything
following it are not asserted.
\item \textbf{Square-free affine algebra.} The assembly
(Proposition~\ref{prop:assembly}) uses the direct-product structure of
\(A\); non-reduced alphabets (e.g.\ chain rings) are outside the scope.
\item \textbf{Moment model.} The means and variances of
Section~\ref{sec:moments} refer to the uniform distribution on labeled
assignments. Orbit contributions are independent in this model; edges
within an orbit are not (Lemma~\ref{lem:varpieces}), and the
\(L=1\), \(L=2\), \(L\ge3\) regimes are genuinely distinct.
\item \textbf{Convention dependence.} The predecessor index
\(\alpha_{i-1}\) in Theorems~\ref{thm:T3}--\ref{thm:T6} is tied to the
code-first pairing of Section~\ref{sec:prelim}; reading the pairing in
the opposite order replaces \(\rho\) by \(\sigma\) and
\(\alpha_{i-1}\) by \(\alpha_{i+1}\). Since the joint distribution is
unchanged under orbit reversal, this convention must be stated, as
here, rather than inferred from the enumerator.
\end{itemize}

\begin{remark}[conditional quantum remark]
In the finite-field case (\(N=1\), \(m_1=1\)), let \(C\) be a code with
parameters \([n, K]\) over \(\Fq\) and let \(H = \dim \Hull_k(C)\). Let
\(\Pi\) be a parity-check matrix of \(C\) (rows a basis of \(C^\perp\)) and
put \(c = \mathrm{rank}\big(\Pi\, \rho(\Pi)^{\mathsf T}\big)\), with
\(\rho\) applied entrywise. The kernel of the pairing map
\(D(C) \to (C^\perp)^*\), \(y \mapsto \big(z \mapsto \sum_i z_i y_i\big)\),
is exactly \(C \cap D(C)\); hence \(c = (n-K) - H\). Therefore, when the
entanglement-assisted construction is used in its Galois form
\cite{DIMP2026} (the framework of \cite{BDS2006}), \(C\) yields an
\([[n, K-H, d_{\mathrm{EA}};\, n-K-H]]\) entanglement-assisted quantum
code, where \(d_{\mathrm{EA}}\) is the distance furnished by that
construction and must be determined independently. The enumerator of
Theorem~\ref{thm:T6} determines the joint distribution of \((K, H)\)
only; it yields no distance information, and no quantum parameter is
claimed from it alone. The affine case requires, in addition, control of
distances under the product decomposition.
\end{remark}

% =====================================================================
\section{Conclusion}
% =====================================================================

For repeated-root constacyclic codes over a square-free affine algebra
with compatible twists, the joint distribution of code dimension and
Galois hull dimension over all labeled exponent assignments is given by
a product of transfer-matrix traces indexed by the orbits of the
reciprocal involution. The enumerator yields the total count, the
code-dimension distribution, the LCD count \(2^{\#\mathrm{orbits}}\), the
self-orthogonal and dual-containing counts, and exact mean and variance
formulas in which length-one orbits enter through the separate constant
\(\nu_P\). At \(P=1\) the theory reduces to the binary simple-root
enumerator. Natural directions for further work include the
incompatible-twist regime, where the dual already exists in a
two-modulus form but the hull intersection must be reconsidered;
counting up to monomial equivalence rather than by labeled assignments;
and distance questions, which the enumerator alone does not address.

% =====================================================================
\begin{thebibliography}{99}
% =====================================================================

\bibitem{BDS2006}
T.~Brun, I.~Devetak, and M.-H.~Hsieh,
Correcting quantum errors with entanglement,
\emph{Science} 314 (2006), 436--439. doi:10.1126/science.1131563.

\bibitem{CDL2015}
B.~Chen, H.~Q. Dinh, and H.~Liu,
Repeated-root constacyclic codes of length \(2\ell^m p^n\),
\emph{Finite Fields and Their Applications} 33 (2015), 137--159.
doi:10.1016/j.ffa.2014.11.006.

\bibitem{DP2025}
I.~Debnath and O.~Prakash,
Average dimensions of Galois hulls of constacyclic codes,
\emph{Advances in Mathematics of Communications} 19(6) (2025),
1569--1604. doi:10.3934/amc.2025010.

\bibitem{DIMP2026}
I.~Debnath, H.~Islam, E.~Mart\'inez-Moro, and O.~Prakash,
Galois hulls of constacyclic codes over affine algebra rings,
\emph{Discrete Mathematics} 349(2) (2026), Article 114750.
doi:10.1016/j.disc.2025.114750; arXiv:2412.08512.

\bibitem{DIYP2026}
I.~Debnath, H.~Islam, S.~Yadav, and O.~Prakash,
Study of small Galois hull dimensions of constacyclic codes,
\emph{Advances in Mathematics of Communications} 22 (2026), 1--18.
doi:10.3934/amc.2025054.

\bibitem{DPI2023}
I.~Debnath, O.~Prakash, and H.~Islam,
Galois hulls of constacyclic codes over finite fields,
\emph{Cryptography and Communications} 15(1) (2023), 111--127.
doi:10.1007/s12095-022-00591-6.

\bibitem{DPI2023corr}
I.~Debnath, O.~Prakash, and H.~Islam,
Correction to: Galois hulls of constacyclic codes over finite fields,
\emph{Cryptography and Communications} 15 (2023), 129--130.
doi:10.1007/s12095-022-00602-6.

\bibitem{ES2026}
R.~F.~Taki Eldin and P.~Sol\'e,
Generator polynomial matrices of the Galois hulls of multi-twisted codes,
\emph{Finite Fields and Their Applications} 109 (2026), Article 102712.
doi:10.1016/j.ffa.2025.102712; arXiv:2305.15237.

\bibitem{FZ2017}
Y.~Fan and L.~Zhang,
Galois self-dual constacyclic codes,
\emph{Designs, Codes and Cryptography} 84 (2017), 473--492.
doi:10.1007/s10623-016-0282-8.

\bibitem{HP2003}
W.~C. Huffman and V.~Pless,
\emph{Fundamentals of Error-Correcting Codes},
Cambridge University Press, 2003.

\bibitem{MS1977}
F.~J. MacWilliams and N.~J.~A. Sloane,
\emph{The Theory of Error-Correcting Codes},
North-Holland, 1977.

\bibitem{Sharma2015}
A.~Sharma,
Repeated-root constacyclic codes of length \(\ell^t p^s\) and their
dual codes,
\emph{Cryptography and Communications} 7 (2015), 229--255.
doi:10.1007/s12095-014-0106-5.

\bibitem{Talbi2021}
S.~Talbi, A.~Batoul, A.~Fotue Tabue, and E.~Mart\'inez-Moro,
Galois hulls of cyclic serial codes over a finite chain ring,
arXiv:2102.06995 (2021).

\bibitem{WZLL2023}
T.~Wu, S.~Zhu, L.~Liu, and L.~Li,
Repeated-root constacyclic codes of length \(6\ell^m p^n\),
\emph{Advances in Mathematics of Communications} 17(5) (2023),
1154--1180. doi:10.3934/amc.2021044.

\bibitem{ZKZ2026}
E.~Zhang, B.~Kong, and X.~Zheng,
Quantum codes from Galois hulls of constacyclic codes over a finite
non-chain ring,
\emph{Entropy} 28 (2026), 407. doi:10.3390/e28040407.

\bibitem{ZWLL2025}
Q.~Zhang, W.~Wang, S.~Luo, and Y.~Li,
Repeated-root constacyclic codes of length \(k^s\ell^m p^n\) over
finite fields,
\emph{Finite Fields and Their Applications} 101 (2025), Article 102542.
doi:10.1016/j.ffa.2024.102542.

\end{thebibliography}

\end{document}
